\section{Conclusion and Future Work}
\label{sec:conclusion}

This paper presented an intent-driven orchestration system for 5G RAN and core network slices using free5GC, OCUDU, Kubernetes, and Nephio. The E2E orchestrator translates service expectations into aligned core and RAN parameters, while the OCUDU and free5GC Operators provide the RAN NSSMF and Core NSSMF/VNFM reconciliation paths. Dynamic core subscriber state is injected through a complementary free5GC-facing interface. Experiments show the intended eMBB/URLLC traffic differentiation, while the measured operator stage confirms that domain reconciliation is not the dominant source of the approximately 30-second management latency; GitOps synchronization is.

\subsection{Limitations}

The prototype demonstrates a control pattern rather than a production-ready slice-management platform. It focuses on two representative service classes, eMBB and URLLC; a production system would need a richer catalog of service profiles, tenant quotas, authentication and authorization policies, and explicit rules for deciding whether a requested intent can be accepted.

The evaluation uses emulated UEs and a virtualized Kubernetes environment, so the throughput results should not be interpreted as a production performance benchmark. They validate configuration effects rather than the maximum capacity of the underlying RAN or core implementation. CPU scheduling, virtual networking, and buffering can influence measured loss, jitter, and throughput. The most important evidence is therefore relative: when the generated slice configuration changes, traffic allocation changes in the expected direction.

\subsection{Future Work}

Future work should add 3GPP Release 18 and Release 19 style intent features, including fulfillment feasibility checks, conflict reports, negotiation, and quantitative utility functions. An admission controller could reject impossible or conflicting intents before they trigger package mutation. For example, two tenants could request full PRB access for different URLLC slices at the same time; the current prototype can detect unsupported input, but it does not yet negotiate alternatives or compute a utility-based compromise.

The hybrid southbound design should also be refined. RAN configuration and free5GC NF deployment already use Nephio/GitOps and domain operators. Dynamic core subscriber and QoS configuration still uses free5GC-facing logic because that state is not fully represented by the operator's deployment resources. A future implementation should migrate more subscriber state into declarative resources so that both domains can share the same package lifecycle, audit trail, and rollback behavior.

On the data plane, deeper integration with Linux traffic control, UPF-level QoS enforcement, or additional OCUDU scheduling hooks could improve end-to-end differentiation beyond RAN scheduling alone. On the management plane, reducing GitOps synchronization latency or allowing policy-controlled fast paths would make the system applicable to a broader range of operational workflows.
